\begin{figure}[htbp]
  \centering
  \begin{tikzpicture}[
    >=Latex,
    font=\small,
    node distance=12mm and 12mm,
    class/.style={
      draw,
      rounded corners=3pt,
      align=left,
      fill=gray!8,
      minimum width=36mm
    },
    arrow/.style={-{Latex[length=2.3mm]}, thick}
  ]
    \node[class] (note) {
      \textbf{Note}\\
      degree\\
      key\\
      midiNum\\
      duration
    };

    \node[class, right=of note] (motif) {
      \textbf{Motif}\\
      beats\\
      notes: Note[]
    };

    \node[class, right=of motif] (bar) {
      \textbf{Bar}\\
      beats\\
      notes: Note[]
    };

    \node[class, below=of motif] (phrase) {
      \textbf{Phrase}\\
      bars: Bar[]\\
      key
    };

    \draw[arrow] (note) -- node[above] {reused in} (motif);
    \draw[arrow] (motif) -- node[above] {fills or shapes} (bar);
    \draw[arrow] (bar) -- node[right] {contained in} (phrase);
  \end{tikzpicture}
  \caption{Simplified object model of the second iteration. The important change is the introduction of explicit musical groupings such as motifs, bars, and phrases, which makes repetition and phrase-level planning possible.}
  \label{fig:iter2_object_model}
\end{figure}
